% ============================================================================
% РЕАЛЬНАЯ ВЕРСИЯ — типографируется при \vtwotrue в preamble.tex
%
% ============================ КАК ЗАПОЛНИТЬ ЭТОТ ФАЙЛ ========================
% 1. Закрыть набор испытаний для лестницы на стенде.
% 2. Направить analysis/build_composition.py на файл запусков композиции. Он
%    должен выдавать, в точности как build_paper.py делает для E1:
%       paper/tables/v1_ladder_solve.tex     доля решённых по ступени + контроль
%       paper/tables/v2_ladder_paired.tex    парный Макнемар против контроля
%       paper/tables/v3_ladder_cost.tex      стоимость/пик/вывод на обеих решённых
%       paper/tables/v4_ladder_redundancy.tex избыточный осмотр по ступени
%    и добавлять его макросы в paper/generated/numbers.tex.
% 3. Установить \vtwotrue в preamble.tex.
% 4. Заменить каждый \pending{} ниже предложением, которое подтверждают данные.
%
% ПРАВИЛО: не вводить измеренное число в этот файл вручную. Если число нужно
% в тексте, добавить макрос в билдере и сослаться на него, в точности как
% делают разделы 06-08. Это сохраняет статью пересчитываемой из сырых
% запусков.
% ============================================================================

\section{Композиция состояния: что должно содержать состояние рантайма}
\label{sec:composition}

\subsection{Мотивация из измеренного отказа}

\pending{один абзац: неуспешное состояние пусто ровно от тех фактов,
которые заново выводятся (\S\ref{sec:mechanism}), а оригинал приписывает
68\% отказов моделей с открытыми весами перезаписи моделью собственного
состояния. Оба указывают в одну сторону: вывести объективную половину
состояния из-под контроля модели.}

\subsection{Лестница}

\pending{шестиступенчатая таблица из заготовки, без изменений — это дизайн,
а не результат}

\subsection{Результаты}

\begin{table}[t]
\centering
\caption{Доля решённых задач по ступени композиции.}
\label{tab:ladder-solve}
\begin{tabular}{@{}lrrr@{}}
\toprule
Условие & Решено / запусков & Доля & 95\% ДИ \\
\midrule
\input{tables/v1_ladder_solve}
\bottomrule
\end{tabular}
\end{table}

\begin{table}[t]
\centering
\caption{Парное сравнение с контрольной группой, точный тест Макнемара.}
\label{tab:ladder-paired}
\begin{tabular}{@{}lrrrrr@{}}
\toprule
Условие & Пар & Только контроль & Только рука & Обе & $p$ \\
\midrule
\input{tables/v2_ladder_paired}
\bottomrule
\end{tabular}
\end{table}

\begin{table}[t]
\centering
\caption{Отношения ресурсов по парам, где решили обе руки.}
\label{tab:ladder-cost}
\begin{tabular}{@{}lrrrrr@{}}
\toprule
Условие & Пар & Стоимость & Пиковый промпт & Вывод & $p$ \\
\midrule
\input{tables/v3_ladder_cost}
\bottomrule
\end{tabular}
\end{table}

\pending{чтение. Три вопроса, на которые нужно ответить явно, в этом
порядке: (a) достигает ли какая-либо ступень способности контрольной
группы, (b) какой единственный компонент сдвигает число сильнее всего,
(c) сохраняется ли преимущество по стоимости на этой ступени.}

\subsection{Проверка механизма: действительно ли упала избыточность?}

\begin{table}[t]
\centering
\caption{Доля избыточного осмотра по ступени, против измеренных \num{0.20}
у контрольной группы и \num{0.58} у только-состояния.}
\label{tab:ladder-redundancy}
\begin{tabular}{@{}lrrr@{}}
\toprule
Условие & Команд / запуск & Доля осмотра & Повторный осмотр \\
\midrule
\input{tables/v4_ladder_redundancy}
\bottomrule
\end{tabular}
\end{table}

\pending{Это гейт, решающий, верна ли история этой статьи. Если ступень
повышает долю решённых задач, но оставляет долю избыточного осмотра на
уровне только-состояния, выигрыш пришёл не от механизма, который
утверждает этот раздел, и утверждение нужно отозвать, а не переформулировать.}

\subsection{Смешение с токенами, сформулированное до того, как это сделает читатель}

При равном $\tailk$ более богатая ступень видит строго больше контекста, чем
более бедная: добавленный компонент — это текст, которого более бедная
ступень не получает. Преимущество одной ступени поэтому частично — это
просто больше входных данных, и его нельзя приводить как «этот компонент
лучше».

\pending{защитимая форма — это утверждение о ФОРМЕ лестницы: постоянное
добавление токенов сдвигает кривую, но не изгибает её. Приводить размеры
промпта и стоимость в той же таблице, а не в сноске, как делает
Таблица~\ref{tab:ladder-cost}.}

\subsection{Намеренно не варьируемое здесь}

\emph{Govern} (перехват предложенной команды и возврат более раннего
результата вместо её выполнения) исключён. Собственная абляция \Ledger{}
приписывает govern основную часть выигрыша по разрешению\origin, так что
его включение примешало бы крупный эффект к измерению меньшего. Это
отдельное условие со своим собственным дизайном.
\note{если дальнейшие гибриды, которые запускает владелец, включают руку,
похожую на govern, ей место в собственном подразделе здесь, а не в
лестнице}
